\section{Conclusion}
\label{sec:conlusion}
ProvSQL's data model, architecture, and generic design makes it suitable
to compute a large range of provenance annotations
and to
provide a generic solution for query evaluation in probabilistic databases. ProvSQL supports
a larger subset of SQL than other existing systems and is actively
maintained. Its performance is suitable for multi-GB
databases; overhead of provenance computation is often
constant (with the notable exception of aggregate queries) and, despite
the \textsf{\#P}-hardness of the problem, it is often possible to compute
probabilities of query outputs in acceptable time.

Many perspectives for improvement exist. First, augmenting the support of
SQL is a necessity for general applicability; this includes
relatively simple cases such as subqueries in \mintkeyword{WHERE}
clauses or \mintkeyword{OUTER JOIN} queries, nested aggregations, and much more complex cases
such as \mintkeyword{WITH RECURSIVE} queries which requires changes to
the
query executor.
Second, when
evaluating safe queries over tuple-independent databases, the intensional
approach to probabilistic query evaluation is suboptimal. It is
unknown whether the safe plan algorithm~\cite{dalvi2013dichotomy}
can be turned into an intensional approach, but some results in this
direction~\cite{DBLP:conf/pods/Monet20} are a first step toward
safe query plans.
